
\subsection{Sensitivity to unobserved confounding}
\label{sec:sensitivity-12}

The Shen-Ross DML estimates in Section~\ref{sec:shen-12city} rest on the
conditional ignorability assumption that, once the structured property covariates
and the location indicators we observe have been partialled out, the residual
variation in listing language is exogenous to the residual variation in log sale
price. Nothing in the data can confirm that assumption directly, so we ask
instead how strong an unobserved confounder would have to be before it overturned
each per-market estimate, following the benchmarking calculus of
\citet{cinelli2020making}, and we pair that bound with a leave-one-out ablation of
the spatial control block, the one observed covariate group most directly
implicated in the text-as-spatial-proxy mechanism that motivates the paper.

\begin{table}[H]
\centering
\caption{Per-market sensitivity of the Shen-Ross estimate to unobserved
confounding, for all twelve markets ordered by robustness value. $RV_{q=1}$ is
the \citet{cinelli2020making} robustness value, the partial $R^2$ an unobserved
confounder must reach on both the treatment and log price to drive the estimate
to zero; $f^2$ is the partial-$R^2$ benchmark of the strongest observed covariate
on the same scale; $\Delta\hat\theta$ is the change in the estimate when the
spatial control block, the latitude-longitude coordinates together with the ZIP
indicators, is dropped. A checkmark marks the single market whose spatial block
carries partial $R^2$ for treatment or outcome exceeding its own robustness
value.}
\label{tab:confounder-sensitivity-12}
\small
\begin{tabular}{lrrrc}
\toprule
Market & $RV_{q=1}$ & $f^2$ & $\Delta\hat\theta_{\mathrm{spatial}}$ & Fragile \\
\midrule
Boston        & $0.005$ & $0.0000$ & $-0.0015$ & \checkmark \\
Seattle       & $0.025$ & $0.0006$ & $+0.0002$ & \\
New York      & $0.030$ & $0.0010$ & $+0.0024$ & \\
Portland      & $0.040$ & $0.0017$ & $-0.0010$ & \\
Washington    & $0.041$ & $0.0017$ & $-0.0011$ & \\
Phoenix       & $0.048$ & $0.0024$ & $+0.0040$ & \\
Dallas        & $0.051$ & $0.0027$ & $-0.0031$ & \\
Atlanta       & $0.051$ & $0.0027$ & $-0.0007$ & \\
Chicago       & $0.064$ & $0.0043$ & $+0.0020$ & \\
Denver        & $0.065$ & $0.0045$ & $+0.0022$ & \\
San Francisco & $0.140$ & $0.0229$ & $-0.0017$ & \\
Philadelphia  & $0.183$ & $0.0411$ & $-0.0035$ & \\
\bottomrule
\end{tabular}
\end{table}

\paragraph{Robustness value.} The robustness value records the minimum strength,
measured as a partial $R^2$ on both the treatment and the outcome, that an
unobserved confounder must reach to drive an estimate to zero, so that a larger
value means a more robust finding. Table~\ref{tab:confounder-sensitivity-12}
reports it for all twelve markets alongside the benchmark strength of the
strongest covariate we actually observe, and Figure~\ref{fig:sensitivity-rv}
orders the markets by it. Two markets clear the conventional $0.10$ threshold,
Philadelphia at $0.183$ and San Francisco at $0.140$, so that overturning either
would require a confounder explaining more than fourteen percent of the residual
variance in both the listing signal and the price. The other ten fall below the
threshold, descending from Denver and Chicago near $0.065$ through the cluster of
mid-sized markets around $0.05$ to New York at $0.030$, Seattle at $0.025$, and
Boston at $0.005$, the last of which has almost no effect to protect.

\paragraph{Observed-covariate benchmark.} The benchmark $f^2$, which places the
most influential observed covariate on the same partial-$R^2$ scale as the
robustness value, never rises above $0.041$, reaching its largest value in
Philadelphia and its smallest, near $2\times10^{-5}$, in Boston. The strongest
covariate we measure therefore sits well inside the robustness value in every
market, most consequentially in Philadelphia and San Francisco where the room
between the two is widest, so that an unobserved confounder capable of overturning
those estimates would have to be considerably stronger than any spatial or
structural feature the data record. The reassurance is sharpest precisely where
the effect is largest, and it thins toward the low-robustness markets, where
neither benchmark is large enough to anchor a confident reading in either
direction. The bias geometry behind these per-market summaries is drawn out, for
the two markets that clear the robustness threshold, in
Figure~\ref{fig:sensitivity-contour}.

\begin{figure}[H]
\centering
\begin{tikzpicture}
\begin{axis}[
    causalre, name=scA, width=0.44\textwidth, height=6.1cm,
    xlabel={Confounder partial $R^2$ with treatment},
    ylabel={Partial $R^2$ with outcome},
    xmin=0, xmax=0.25, ymin=0, ymax=0.25,
    xtick={0,0.05,0.10,0.15,0.20}, ytick={0,0.05,0.10,0.15,0.20},
    title={(a) San Francisco}, title style={font=\small},
    legend style={font=\scriptsize, draw=none, fill=none,
                  at={(0.0,1.15)}, anchor=south west, legend columns=4,
                  /tikz/every even column/.append style={column sep=6pt}},
]
\addplot[gray, dotted, thick, domain=0:0.25, forget plot] {x};
\addplot[cSF, thick, domain=0.085:0.25, samples=90] {0.02292*(1-x)/x};
\addlegendentry{$\hat\theta_{\mathrm{adj}}=0$}
\addplot[cSF, dashed, domain=0.023:0.25, samples=90] {0.00573*(1-x)/x};
\addlegendentry{$\hat\theta_{\mathrm{adj}}=\hat\theta/2$}
\addplot[only marks, mark=*, mark size=2pt, color=black] coordinates {(0.140,0.140)};
\addlegendentry{robustness value}
\addplot[only marks, mark=triangle*, mark size=2.6pt, color=cNYC] coordinates {(0.023,0.023)};
\addlegendentry{strongest covariate}
\end{axis}
\begin{axis}[
    causalre, width=0.44\textwidth, height=6.1cm,
    at={(scA.east)}, anchor=west, xshift=1.0cm,
    xlabel={Confounder partial $R^2$ with treatment}, ylabel={},
    xmin=0, xmax=0.25, ymin=0, ymax=0.25,
    xtick={0,0.05,0.10,0.15,0.20}, ytick={0,0.05,0.10,0.15,0.20},
    title={(b) Philadelphia}, title style={font=\small},
]
\addplot[cSF, thick, domain=0.142:0.25, samples=90] {0.04114*(1-x)/x};
\addplot[cSF, dashed, domain=0.040:0.25, samples=90] {0.01029*(1-x)/x};
\addplot[gray, dotted, thick, domain=0:0.25] {x};
\addplot[only marks, mark=*, mark size=2pt, color=black] coordinates {(0.183,0.183)};
\addplot[only marks, mark=triangle*, mark size=2.6pt, color=cNYC] coordinates {(0.041,0.041)};
\end{axis}
\end{tikzpicture}
\caption{Omitted-variable-bias sensitivity contours, in the manner of
\citet{cinelli2020making}, for the two markets that clear the robustness
threshold. The axes are the partial $R^2$ that a single unobserved confounder
would have with the treatment (horizontal) and with the outcome (vertical). The
solid curve is the locus along which such a confounder would drive the
per-standard-deviation estimate to zero, and the dashed curve the locus along
which it would halve it. The black dot, where the solid curve meets the dotted
equal-strength diagonal, is the robustness value, the equal-strength confounder
that just eliminates the effect, at $0.14$ in San Francisco and $0.18$ in
Philadelphia. The triangle marks a confounder as strong
as the most influential covariate actually observed, which falls far inside the
kill curve in both markets, so an unobserved confounder would have to be several
times stronger than any measured one to overturn the estimate.}
\label{fig:sensitivity-contour}
\end{figure}

\paragraph{Spatial-block leave-one-out.} Dropping the spatial control block and
re-estimating moves the point estimate by at most four thousandths of a log point
in any of the twelve markets, the largest shift falling in Phoenix at $+0.004$
and the rest smaller still. That the estimates barely respond to removing the
location controls is the pattern one would hope to see if the text signal were not
merely spatial information re-expressed, since a coefficient that rode on
geography would collapse once geography left the conditioning set rather than hold
to within Monte Carlo noise of where it began. The leave-one-out thus localizes,
market by market, the conclusion that the deconfounding decomposition of
Section~\ref{sec:decomposition} reaches for the embedding channel as a whole.

\paragraph{Fragility.} Under the same benchmarking convention a market is flagged
when the partial $R^2$ that the spatial block carries for either the treatment or
the outcome exceeds that market's robustness value, the condition under which a
confounder no stronger than the location controls already in the model could in
principle account for the estimate. Only Boston trips the flag, and it does so
because its robustness value is itself almost zero, the spatial block's
explanatory power for the treatment, at $0.012$, edging past a robustness value of
$0.005$ for an estimate that is statistically indistinguishable from zero to begin
with. No market with a robustness value above $0.025$ is flagged, so the fragility
check isolates the one place where there is no effect to confound rather than
casting doubt on any market where an effect is present.

\paragraph{Interpretation.} The analysis is deliberately confined to the single
observed covariate group the per-market panels share and the one most tied to the
mechanism this paper studies, so the robustness value carries the general bound
while the spatial leave-one-out localizes how much of each estimate rides on
measured location. The two diagnostics agree with each other and with the rest of
the section. Philadelphia and San Francisco, the markets in which the text effect
is largest, are also the ones most robust to unobserved confounding; the markets
that fall below the threshold are those whose estimates are already small; and
removing the explicit location controls disturbs none of them. We read the
per-market sensitivities as informative about robustness rather than as a second
identification test, and we leave the benchmarking of non-spatial confounder
groups to future work, since the census, crime, and amenity blocks needed to
construct them are not uniformly available across the twelve per-market panels.

\begin{figure}[H]
\centering
\begin{tikzpicture}
\begin{axis}[
    causalre, width=0.62\textwidth, height=6.6cm,
    symbolic y coords={Boston,Seattle,New York,Portland,Washington,Phoenix,Dallas,Atlanta,Chicago,Denver,San Francisco,Philadelphia},
    ytick=data, enlarge y limits=0.06,
    xlabel={Robustness value $RV_{q=1}$ (residual share an unmeasured confounder must explain in both treatment and price)},
    xmin=0, xmax=0.21, xtick={0,0.05,0.10,0.15,0.20},
    xticklabel={\pgfmathprintnumber[fixed,precision=2]{\tick}},
    title={Sensitivity of the per-market estimate to unmeasured confounding},
]
\addplot[only marks, mark=*, mark size=2.0pt, color=cNYC] coordinates {
  (0.005,Boston) (0.025,Seattle) (0.030,New York) (0.040,Portland) (0.041,Washington)
  (0.048,Phoenix) (0.051,Dallas) (0.051,Atlanta) (0.064,Chicago) (0.065,Denver)
  (0.140,San Francisco) (0.183,Philadelphia)};
\end{axis}
\end{tikzpicture}
\caption{Robustness value for each market's sentence-BERT effect, the share of
residual variance an unmeasured confounder would need to explain in both the
treatment and log price to drive the estimate to zero, following
\citet{cinelli2020making}. Larger values are more robust. The estimate is most
robust in Philadelphia ($0.18$) and San Francisco ($0.14$) and most fragile in
the markets whose point estimates are already near zero. The modest robustness
values in several markets are the reason we read the per-market estimates as
associational and sensitivity-bounded rather than cleanly causal.}
\label{fig:sensitivity-rv}
\end{figure}

\subsection{A composition-and-conditioning audit for text-treatment DML}
\label{sec:coca}

The robustness values of Table~\ref{tab:confounder-sensitivity-12} benchmark the
strength an unobserved confounder would need against the strength of the strongest
confounder we do observe, following \citet{cinelli2020making}. That benchmark has a
blind spot that a text-derived treatment walks directly into, and this section
names it, bounds it, and turns it into a pre-registered diagnostic.

Consider an omitted categorical covariate $C$, property type being the concrete
instance and vacant land the sharpest case, with three properties that hold in the
data. It shifts the treatment, $R^2_{T \sim C \mid X} > 0$, because a vacant lot is
described in different language than a house. It shifts price net of the treatment
and the observed controls, $R^2_{Y \sim C \mid T, X} > 0$. And it is nearly
invisible to the observed confounder block, $R^2(C \mid X) \approx 0$: a
cross-fitted classifier of the vacant-land indicator on the full confounder set
runs at an AUC of $0.50$ in Boston and $0.47$ in New York, while the same indicator
is recovered from the text direction at an AUC above $0.95$. The nearest-parcel
join that supplies the structured controls cannot represent landness, since for a
vacant lot it returns a neighbouring built parcel's attributes, so the design is
blind to exactly the attribute the text encodes. The bias from omitting such a $C$
is bounded by the partial-$R^2$ omitted-variable-bias inequality of
\citet{cinelli2020making} and its double-machine-learning generalization in
\citet{chernozhukov2022long}, and that bound holds whether or not $C$ is orthogonal
to $X$. What the orthogonality changes is not the bound but the benchmark used to
fill it in.

\begin{proposition}[Orthogonal composition defeats the observed-covariate benchmark]
\label{prop:coca}
In the partially-linear model of Equation~\eqref{eq:plr}, let $C$ be an omitted
covariate with $r_D := R^2_{T \sim C \mid X} > 0$ and
$r_Y := R^2_{Y \sim C \mid T, X} > 0$, and suppose $R^2(C \mid X) \to 0$. Then (i)
the protective subtraction that conditioning on $X$ applies to a confounder
correlated with the controls, the cross term $\rho_{TX}\rho_{CX}$ in the partial
correlation, vanishes, so the partial strengths $r_D$ and $r_Y$ that enter the
\citet{cinelli2020making} and \citet{chernozhukov2022long} bias bound are not
shrunk below their raw values, with $r_D \to R^2_{T \sim C}/(1 - R^2_{T \sim X})$ in
the limit; and (ii) the
observed-covariate benchmark, which calibrates the plausible magnitude of $C$ by
that of the strongest observed covariate, carries no warrant, because $C$ by
construction carries information orthogonal to the covariate space on which that
benchmark is computed. The robustness value reported against the observed benchmark
is therefore numerically well defined but uninformative about $C$.
\end{proposition}

The argument is short. Partialling $X$ out of $T$ and of $C$ leaves residuals whose
correlation obeys the recursive identity
$\rho_{TC \mid X} = (\rho_{TC} - \rho_{TX}\rho_{CX}) /
\sqrt{(1 - \rho_{TX}^2)(1 - \rho_{CX}^2)}$, whose protective subtraction
$\rho_{TX}\rho_{CX}$ vanishes as $\mathrm{Cor}(C, X) \to 0$, leaving
$\rho_{TC \mid X} \to \rho_{TC} / \sqrt{1 - \rho_{TX}^2}$, no smaller than the raw
correlation; the same holds on the outcome side, which gives part (i). A computer-
algebra check confirms both this identity and the fact that for a single omitted
confounder the \citet{cinelli2020making} bias expression equals the exact bias
rather than merely bounding it, so the bound the audit reasons about is tight; the
verification is released with the replication code. Part (ii)
is the operational reading: the benchmark's premise is that an omitted confounder
is comparable in strength to an observed covariate, and a classifier of $C$ on $X$
running at chance is direct evidence that $C$ lives outside the space that premise
is calibrated on. We are careful about what this does and does not establish. It is
not that the formal bound cannot detect $C$; \citet{cinelli2020making} show in
their own Section~6.2 that the correctly computed bound recovers the right answer
for a confounder orthogonal to the controls when the analyst sets the benchmark
strength correctly. It is that the routine heuristic of setting that strength by
analogy to observed covariates has no anchor once $R^2(C \mid X) \approx 0$, and
this is precisely the regime in which an analyst is most likely to feel wrongly
reassured, since orthogonal to the controls reads too easily as weak. The claim in
part (i) that conditioning shrinks a correlated confounder's residual strength is
generic rather than universal, since \citet{middleton2016bias} show that
conditioning on covariates correlated with an unobserved confounder can amplify
rather than shrink the bias in bad-control configurations, so it is the expected
case and not a law.

A second, distinct failure channel travels with the first. When $C$ nearly sorts
the treatment, the residual treatment variance $E[(T - E[T \mid X, C])^2]$
collapses and the orthogonal score becomes ill conditioned, a phenomenon
\citet{saco2025illconditioned} characterize through the conditioning number
$\kappa = 1/(1 - R^2_{\text{oof}}(T \mid X))$; that reference reports the diagnostic
as a fragility assessment and deliberately proposes no universal threshold, and we
follow it in treating $\kappa$ and the design-matrix condition number as flags read
alongside the substantive analysis rather than as a test. The composition confound
and the conditioning degeneracy are separate objects, and confusing one for the
other misdiagnoses the fix: the first is corrected by adjusting for the omitted
category, the second by repairing the design.

These two channels motivate a two-step audit, run before the point estimate is
inspected, that we call the composition-and-conditioning audit. Step one screens
conditioning: flag a market if the standardized confounder matrix has condition
number above $10^6$ or any confounder column has collapsed to a constant, either of
which signals a design-matrix degeneracy rather than a treatment effect. Step two
screens composition: flag a market if a candidate omitted category is
chance-predictable from the controls, $\mathrm{AUC}(C \mid X) < 0.60$, yet strongly
predictable from the treatment, $\mathrm{AUC}(C \mid T) > 0.75$. The prescribed
correction differs by step. A conditioning flag is a pipeline repair. A composition
flag calls for adjusting for the category, and the natural instinct, adding a
property-type dummy to the regularized nuisance, backfires, because the ridge
penalty shrinks a rare category toward no effect: in Boston, where non-residential
listings are $1.9\%$ of the sample, adding the dummy raises $|\hat\theta|$ from
$0.40$ to $0.45$. Exact within-category demeaning of the treatment and the outcome,
which is unpenalized, is the correct adjustment and lowers the Boston estimate to
$0.34$. Across the panel the composition adjustment materially moves only the two
markets with the largest non-residential share, Philadelphia and Chicago, and
leaves the rest, including the New York exception, essentially unchanged, which is
consistent with the mechanism the proposition describes.

A calibration study confirms that the two steps discriminate the two failure modes
rather than firing together. Across a grid of category shares and confounder
strengths, with $200$ replications per cell, the conditioning step fires on a
simulated design degeneracy in every replication and never on a composition
confound, while the composition step fires on a material composition confound, one
strong enough to lift $\mathrm{AUC}(C \mid T)$ above the threshold, in more than
nine of ten replications and never on a design degeneracy; the false-positive rate
of each step on the other's failure mode is zero. The audit therefore separates a
numerical artifact from a substantive confound on data the analysis already has,
and it does so before the estimate is read, which is what a pre-registered
diagnostic is meant to do. We claim for it neither a new estimator nor a new limit
theorem, only the connection of two existing literatures, the omitted-variable-bias
bounds of \citet{cinelli2020making} and \citet{chernozhukov2022long} and the
conditioning diagnostics of \citet{saco2025illconditioned}, into a screen fitted to
the specific way a text-derived treatment fails.
